\section{核验锚点簇：正德、嘉靖：朝廷程序、海防与地方压力（1506--1567）}
\label{register:1506-1567_zhengde_jiajing}
宫廷程序、东南海岸与北方边镇在此并列，不表示它们由一项政策驱动。记录只把可见的诏令、调兵、港口和官署接回各自的证据链，执行范围仍须逐地核验。

以下锚点只承担年序导航；本章前文负责叙事。完整来源、校勘与材料限度留在来源附录和必要的信息框，不把账本字段重复铺成读者段落。

\begin{description}
  \item[\eventanchor{ML-1506-001}{1506（正德元年）五月：税务部责令调查收入不足问题}]
  \item[\eventanchor{ML-1506-002}{1506（正德元年）七月：户部尚书韩文抱怨皇帝用大臣府库支出}]
  \item[\eventanchor{ML-1506-003}{1506（正德元年）十月28日：户部尚书奏请皇帝将他亲手的太监全部处死，皇帝不肯，大臣们都辞职了。}]
  \item[\eventanchor{ML-1506-004}{1506（正德元年）：正德皇帝乔装游走北京街头}]
  \item[\eventanchor{ML-1507-001}{1507（正德二年）九月：元宵节花灯费35万两银子}]
  \item[\eventanchor{ML-1507-002}{1507（正德二年）：科举、讲学和官员文集的本年线索可用于追踪知识网络，主张与行政结果须分开。}]
  \item[\eventanchor{ML-1507-003}{1507（正德二年）：嘉靖初继统礼制的本年文书构成大礼议过程锚点，礼制理由和政治效应不可简化为单一动机。}]
  \item[\eventanchor{MM-1507-LIU-JIN-FISCAL-EXTRACTIONS}{1507（正德二年）：刘瑾集团要求地方严核钱粮、差役与库藏，内廷对财政文书的介入加深}]
  \item[\eventanchor{MM-1507-LIU-JIN-PURGE}{1507（正德二年）：刘瑾借厂卫和考察处置异己官员，朝廷人事与言路受到冲击}]
  \item[\eventanchor{ML-1508-002}{1508（正德三年）：科举、讲学和官员文集的本年线索可用于追踪知识网络，主张与行政结果须分开。}]
  \item[\eventanchor{ML-1508-003}{1508（正德三年）：嘉靖初继统礼制的本年文书构成大礼议过程锚点，礼制理由和政治效应不可简化为单一动机。}]
  \item[\eventanchor{ML-1508-004}{1508（正德三年）：正德朝内廷、人事与巡幸相关文书构成本年中枢政治线索，后出道德化叙事须与原奏疏分开。（材料核查重点：报称信息的来源、抄传与行政分类）}]
  \item[\eventanchor{MM-1508-LIU-JIN-LAND-SURVEY}{1508（正德三年）：刘瑾推动清丈田亩和赋役整饬，以核查隐漏田土}]
  \item[\eventanchor{ML-1509-001}{1509（正德四年）八月：辽东两驻军叛乱，分白银2500两后平息}]
  \item[\eventanchor{ML-1509-002}{1509（正德四年）：地方反抗、边镇调兵和宗藩危机的本年军政材料标记跨区域动员，军数与战果需要复核。}]
  \item[\eventanchor{ML-1509-003}{1509（正德四年）：嘉靖初继统礼制的本年文书构成大礼议过程锚点，礼制理由和政治效应不可简化为单一动机。}]
  \item[\eventanchor{ML-1509-004}{1509（正德四年）：正德朝内廷、人事与巡幸相关文书构成本年中枢政治线索，后出道德化叙事须与原奏疏分开。}]
  \item[\eventanchor{ML-1510-001}{1510（正德五年）五月12日：安化王叛乱：山西朱之璠叛乱}]
  \item[\eventanchor{ML-1510-002}{1510（正德五年）五月30日：安化王叛乱：朱之璠被俘}]
  \item[\eventanchor{ML-1510-003}{1510（正德五年）：达延汗征服鄂尔多斯河套地区}]
  \item[\eventanchor{MM-1510-ANHUA-REBELLION}{1510（正德五年）五月：安化王朱寘鐇在宁夏起兵，旋被当地官军与将领平定}]
  \item[\eventanchor{MM-1510-LIU-JIN-EXECUTION}{1510（正德五年）八月：刘瑾被以谋反等罪处死}]
  \item[\eventanchor{ML-1511-001}{1511（正德六年）二月：北京周边土匪起义}]
  \item[\eventanchor{ML-1511-002}{1511（正德六年）十月：北京各地土匪烧毁皇家粮车}]
  \item[\eventanchor{ML-1511-003}{1511（正德六年）：占领马六甲 (1511)：马六甲苏丹国向葡萄牙人发出求援请求}]
  \item[\eventanchor{MM-1511-LIU-LIU-QI-UPRISING}{1511（正德六年）：刘六、刘七等在山东、北直隶一带聚众活动，成为跨区域治安危机}]
  \item[\eventanchor{ML-1512-002}{1512（正德七年）九月7日：强盗大军被击败}]
  \item[\eventanchor{ML-1512-003}{1512（正德七年）：地方反抗、边镇调兵和宗藩危机的本年军政材料标记跨区域动员，军数与战果需要复核。}]
  \item[\eventanchor{MM-1512-LIU-LIU-QI-SUPPRESSION}{1512（正德七年）：官军在山东、北直隶分路追击刘六、刘七等集团，主要武装至年内瓦解}]
  \item[\eventanchor{ML-1513-001}{1513（正德八年）：清丈、库藏、差役与征解的本年记录显示财政监管压力，整饬命令不等于隐漏已被清除。（材料核查重点：诏令或奏议的形成与发受链）}]
  \item[\eventanchor{ML-1513-002}{1513（正德八年）：正德朝内廷、人事与巡幸相关文书构成本年中枢政治线索，后出道德化叙事须与原奏疏分开。}]
  \item[\eventanchor{ML-1513-003}{1513（正德八年）：清丈、库藏、差役与征解的本年记录显示财政监管压力，整饬命令不等于隐漏已被清除。（材料核查重点：报称信息的来源、抄传与行政分类）}]
  \item[\eventanchor{MM-1513-PORTUGUESE-TAMAO-CONTACT}{1513（正德八年）：葡萄牙航海者在屯门一带停泊并尝试贸易，为后续使团接触提供海上先例}]
  \item[\eventanchor{ML-1514-001}{1514（正德九年）二月10日：宫殿庭院内的火药帐篷着火，烧毁了住宅宫殿}]
  \item[\eventanchor{ML-1514-002}{1514（正德九年）九月：正德皇帝被老虎咬成重伤}]
  \item[\eventanchor{ML-1514-003}{1514（正德九年）：东南港口、日本使团或葡萄牙接触的本年材料为海上交往留档，朝贡、私贸与冲突不可混称。}]
  \item[\eventanchor{MM-1514-TURPAN-HAMI-ANNEXATION}{1514（正德九年）：吐鲁番势力再次夺取哈密，明朝维系该绿洲朝贡政权的能力显著下降}]
  \item[\eventanchor{ML-1515-001}{1515（正德十年）夏：京师、侍卫三万出动，重建宫殿}]
  \item[\eventanchor{ML-1515-002}{1515（正德十年）：东南港口、日本使团或葡萄牙接触的本年材料为海上交往留档，朝贡、私贸与冲突不可混称。}]
  \item[\eventanchor{ML-1515-003}{1515（正德十年）：地方反抗、边镇调兵和宗藩危机的本年军政材料标记跨区域动员，军数与战果需要复核。}]
  \item[\eventanchor{ML-1515-004}{1515（正德十年）：科举、讲学和官员文集的本年线索可用于追踪知识网络，主张与行政结果须分开。}]
  \item[\eventanchor{ML-1516-001}{1516（正德十一年）：地方反抗、边镇调兵和宗藩危机的本年军政材料标记跨区域动员，军数与战果需要复核。}]
  \item[\eventanchor{ML-1516-002}{1516（正德十一年）：科举、讲学和官员文集的本年线索可用于追踪知识网络，主张与行政结果须分开。}]
  \item[\eventanchor{ML-1516-003}{1516（正德十一年）：清丈、库藏、差役与征解的本年记录显示财政监管压力，整饬命令不等于隐漏已被清除。}]
  \item[\eventanchor{ML-1517-001}{1517（正德十二年）十月16日：达延汗袭击明朝}]
  \item[\eventanchor{ML-1517-002}{1517（正德十二年）十月20日：正德皇帝击退达延汗的袭击}]
  \item[\eventanchor{ML-1517-003}{1517（正德十二年）：托梅·皮雷抵达广州}]
  \item[\eventanchor{MM-1517-PORTUGUESE-EMBASSY}{1517（正德十二年）：葡萄牙使团抵达广东，托名满剌加使者请求入朝}]
  \item[\eventanchor{ML-1518-001}{1518（正德十三年）一月：正德皇帝因没有给他足够的钱财而被囚禁在北京}]
  \item[\eventanchor{ML-1518-002}{1518（正德十三年）：东南港口、日本使团或葡萄牙接触的本年材料为海上交往留档，朝贡、私贸与冲突不可混称。}]
  \item[\eventanchor{ML-1518-003}{1518（正德十三年）：正德朝内廷、人事与巡幸相关文书构成本年中枢政治线索，后出道德化叙事须与原奏疏分开。}]
  \item[\eventanchor{MM-1518-WANG-YANGMING-GANZHOU}{1518（正德十三年）：王守仁在赣南主持军政，推行保甲、招抚与军事征剿并用}]
  \item[\eventanchor{ML-1519-001}{1519（正德十四年）七月9日：宁王叛乱：江西朱辰灏叛乱}]
  \item[\eventanchor{ML-1519-002}{1519（正德十四年）七月13日：宁王叛乱：叛军攻克九江}]
  \item[\eventanchor{ML-1519-003}{1519（正德十四年）七月23日：宁王叛乱：叛军围攻安庆}]
  \item[\eventanchor{ML-1519-004}{1519（正德十四年）八月9日：宁王叛乱：叛军解围安庆}]
  \item[\eventanchor{ML-1519-005}{1519（正德十四年）八月13日：宁王叛乱：清军攻克南昌}]
  \item[\eventanchor{ML-1519-006}{1519（正德十四年）八月15日：宁王叛乱：朱辰灏军大败}]
  \item[\eventanchor{ML-1519-007}{1519（正德十四年）八月20日：宁王叛乱：朱辰浩逃离舰队被俘}]
  \item[\eventanchor{MM-1519-WUZONG-SOUTHERN-TOUR}{1519（正德十四年）：宁王之变已受控制后，武宗仍南巡并以亲临军务为名驻跸江南}]
  \item[\eventanchor{ML-1520-001}{1520（正德十五年）一月：正德皇帝禁止宰猪}]
  \item[\eventanchor{ML-1520-003}{1520（正德十五年）：科举、讲学和官员文集的本年线索可用于追踪知识网络，主张与行政结果须分开。}]
  \item[\eventanchor{ML-1520-004}{1520（正德十五年）：地方反抗、边镇调兵和宗藩危机的本年军政材料标记跨区域动员，军数与战果需要复核。}]
  \item[\eventanchor{ML-1521-001}{1521（正德十六年）四月20日：正德皇帝驾崩}]
  \item[\eventanchor{ML-1521-002}{1521（正德十六年）四月21日：托梅·皮雷的政党被驱逐出北京}]
  \item[\eventanchor{ML-1521-003}{1521（正德十六年）四月27日：朱厚熜即位嘉靖皇帝}]
  \item[\eventanchor{ML-1521-004}{1521（正德十六年）五月：屯门之战：明军将拒绝离开的葡萄牙舰队驱逐出屯门}]
  \item[\eventanchor{ML-1521-005}{1521（正德十六年）：宫殿重建完成}]
  \item[\eventanchor{MM-1521-JIAJING-ACCESSION}{1521（嘉靖元年）四月：武宗无子而死，兴王朱厚熜入继帝位，次年改元嘉靖}]
  \item[\eventanchor{MM-1521-YANG-TINGHE-REFORMS}{1521（嘉靖元年）：杨廷和等在嘉靖初裁革部分正德末近侍机构和冗费}]
  \item[\eventanchor{ML-1522-001}{1522（嘉靖元年）：禁止葡萄牙人在广州进行贸易}]
  \item[\eventanchor{ML-1522-002}{1522（嘉靖元年）：山草湾海战：葡萄牙舰队在大屿山附近冲破明朝封锁线，伤亡惨重成功撤离}]
  \item[\eventanchor{ML-1522-003}{1522（嘉靖元年）：清丈、库藏、差役与征解的本年记录显示财政监管压力，整饬命令不等于隐漏已被清除。}]
  \item[\eventanchor{ML-1523-001}{1523（嘉靖二年）五月：宁波事件：细川商团攻击大内商团，抢劫宁波，夺取船只，并在起航前杀死一名明将；中国朝贡体系失去海上贸易价值}]
  \item[\eventanchor{ML-1523-002}{1523（嘉靖二年）：明朝根据葡萄牙设计生产后装旋转枪。}]
  \item[\eventanchor{ML-1523-003}{1523（嘉靖二年）：东南港口、日本使团或葡萄牙接触的本年材料为海上交往留档，朝贡、私贸与冲突不可混称。}]
  \item[\eventanchor{MM-1523-NINGBO-TRIBUTE-CONFLICT}{1523（嘉靖二年）：两支日本朝贡使团在宁波因勘合、赏赐与入港次序冲突并发生暴力，明廷随后收紧相关朝贡管理}]
  \item[\eventanchor{ML-1524-001}{1524（嘉靖三年）八月：大同叛军驻军}]
  \item[\eventanchor{ML-1524-002}{1524（嘉靖三年）：明吐鲁番冲突：吐鲁番攻打甘州，被击退}]
  \item[\eventanchor{ML-1525-001}{1525（嘉靖四年）四月：大同起义军被击败}]
  \item[\eventanchor{ML-1525-002}{1525（嘉靖四年）：嘉靖倭口突袭：双榆成为贸易飞地}]
  \item[\eventanchor{ML-1525-003}{1525（嘉靖四年）：一些福建商人会说台湾语}]
  \item[\eventanchor{ML-1525-004}{1525（嘉靖四年）：嘉靖朝礼制、内阁与人事的本年诏令和奏议形成中枢协调线索，留中与施行之间应予区分。}]
  \item[\eventanchor{MM-1525-YANG-TINGHE-DISMISSAL}{1525（嘉靖四年）：杨廷和在大礼议冲突中失势致仕，嘉靖朝内阁权力格局改变}]
  \item[\eventanchor{ML-1526-001}{1526（嘉靖五年）：北京遭遇饥荒}]
  \item[\eventanchor{ML-1526-002}{1526（嘉靖五年）：嘉靖朝礼制、内阁与人事的本年诏令和奏议形成中枢协调线索，留中与施行之间应予区分。}]
  \item[\eventanchor{ML-1526-003}{1526（嘉靖五年）：海禁、朝贡和沿海港口管理的本年材料显示官方海上政策，禁令范围与实际航行须分别调查。（材料核查重点：诏令或奏议的形成与发受链）}]
  \item[\eventanchor{ML-1526-004}{1526（嘉靖五年）：海禁、朝贡和沿海港口管理的本年材料显示官方海上政策，禁令范围与实际航行须分别调查。（材料核查重点：报称信息的来源、抄传与行政分类）}]
  \item[\eventanchor{ML-1526-005}{1526（嘉靖五年）：北边警报、边镇防务和西南处置的本年军政材料标记边疆压力，敌我称谓与军数应回到原文。}]
  \item[\eventanchor{ML-1527-001}{1527（嘉靖六年）：湖广洪水泛滥}]
  \item[\eventanchor{ML-1527-002}{1527（嘉靖六年）：祀典、学校、出版或讲学的本年材料记录文化制度活动，规范话语不等于社会普遍认同。（材料核查重点：诏令或奏议的形成与发受链）}]
  \item[\eventanchor{ML-1527-003}{1527（嘉靖六年）：祀典、学校、出版或讲学的本年材料记录文化制度活动，规范话语不等于社会普遍认同。（材料核查重点：报称信息的来源、抄传与行政分类）}]
  \item[\eventanchor{ML-1527-004}{1527（嘉靖六年）：赋役折色、盐课、河工或田土核查的本年文书为财政改革史提供节点，制度名称不可跨区套用。}]
  \item[\eventanchor{MM-1527-WANG-YANGMING-GUANGXI}{1527（嘉靖六年）：王守仁受命处理思恩、田州土司冲突，兼用招抚、军事与设治方案}]
  \item[\eventanchor{ML-1528-001}{1528（嘉靖七年）：明吐鲁番冲突：吐鲁番恢复贸易特权}]
  \item[\eventanchor{ML-1528-002}{1528（嘉靖七年）：北边警报、边镇防务和西南处置的本年军政材料标记边疆压力，敌我称谓与军数应回到原文。}]
  \item[\eventanchor{ML-1528-004}{1528（嘉靖七年）：灾异、赈济及地方工程的本年报告可与州县材料互校，救济命令不自动证明救济到达。（材料核查重点：报称信息的来源、抄传与行政分类）}]
  \item[\eventanchor{ML-1528-005}{1528（嘉靖七年）：海禁、朝贡和沿海港口管理的本年材料显示官方海上政策，禁令范围与实际航行须分别调查。}]
  \item[\eventanchor{ML-1529-001}{1529（嘉靖八年）：嘉靖倭寇袭扰：温州数名将领因与倭寇勾结被流放}]
  \item[\eventanchor{ML-1529-002}{1529（嘉靖八年）：看到了一颗不吉利的彗星}]
  \item[\eventanchor{ML-1529-003}{1529（嘉靖八年）：嘉靖朝礼制、内阁与人事的本年诏令和奏议形成中枢协调线索，留中与施行之间应予区分。}]
  \item[\eventanchor{ML-1529-004}{1529（嘉靖八年）：祀典、学校、出版或讲学的本年材料记录文化制度活动，规范话语不等于社会普遍认同。}]
  \item[\eventanchor{MM-1529-WANG-YANGMING-DEATH}{1529（嘉靖七年）：王守仁自广西返乡途中病逝，其学术与军政文字随后由门人整理传播}]
  \item[\eventanchor{ML-1530-001}{1530（嘉靖九年）：赋役折色、盐课、河工或田土核查的本年文书为财政改革史提供节点，制度名称不可跨区套用。}]
  \item[\eventanchor{ML-1530-002}{1530（嘉靖九年）：嘉靖朝礼制、内阁与人事的本年诏令和奏议形成中枢协调线索，留中与施行之间应予区分。}]
  \item[\eventanchor{ML-1530-003}{1530（嘉靖九年）：北边警报、边镇防务和西南处置的本年军政材料标记边疆压力，敌我称谓与军数应回到原文。}]
  \item[\eventanchor{ML-1530-004}{1530（嘉靖九年）：灾异、赈济及地方工程的本年报告可与州县材料互校，救济命令不自动证明救济到达。}]
  \item[\eventanchor{MM-1530-SUBURBAN-RITES}{1530（嘉靖九年）：嘉靖调整郊祀，分别祭天与祭地，并改建相应坛制}]
  \item[\eventanchor{ML-1531-001}{1531（嘉靖十年）：大同遭到蒙古人的袭击}]
  \item[\eventanchor{ML-1531-002}{1531（嘉靖十年）：看到了一颗不吉利的彗星}]
  \item[\eventanchor{ML-1531-003}{1531（嘉靖十年）：赋役折色、盐课、河工或田土核查的本年文书为财政改革史提供节点，制度名称不可跨区套用。}]
  \item[\eventanchor{ML-1531-004}{1531（嘉靖十年）：嘉靖朝礼制、内阁与人事的本年诏令和奏议形成中枢协调线索，留中与施行之间应予区分。}]
  \item[\eventanchor{MM-1531-ONE-WHIP-PROPOSAL}{1531（嘉靖十年）：桂萼等讨论合并赋役、折纳银两的方案，显示一条鞭式改革的早期政策脉络}]
  \item[\eventanchor{ML-1532-001}{1532（嘉靖十一年）：嘉靖倭寇大袭：广州巡抚因铲除倭寇不力被罢免}]
  \item[\eventanchor{ML-1532-003}{1532（嘉靖十一年）：北边警报、边镇防务和西南处置的本年军政材料标记边疆压力，敌我称谓与军数应回到原文。（材料核查重点：诏令或奏议的形成与发受链）}]
  \item[\eventanchor{ML-1532-004}{1532（嘉靖十一年）：海禁、朝贡和沿海港口管理的本年材料显示官方海上政策，禁令范围与实际航行须分别调查。}]
  \item[\eventanchor{ML-1532-005}{1532（嘉靖十一年）：北边警报、边镇防务和西南处置的本年军政材料标记边疆压力，敌我称谓与军数应回到原文。（材料核查重点：报称信息的来源、抄传与行政分类）}]
  \item[\eventanchor{ML-1533-001}{1533（嘉靖十二年）十月24日：大同卫军叛乱，被镇压}]
  \item[\eventanchor{ML-1533-002}{1533（嘉靖十二年）：海禁、朝贡和沿海港口管理的本年材料显示官方海上政策，禁令范围与实际航行须分别调查。}]
  \item[\eventanchor{ML-1533-003}{1533（嘉靖十二年）：赋役折色、盐课、河工或田土核查的本年文书为财政改革史提供节点，制度名称不可跨区套用。（材料核查重点：诏令或奏议的形成与发受链）}]
  \item[\eventanchor{ML-1533-004}{1533（嘉靖十二年）：祀典、学校、出版或讲学的本年材料记录文化制度活动，规范话语不等于社会普遍认同。}]
  \item[\eventanchor{ML-1533-005}{1533（嘉靖十二年）：赋役折色、盐课、河工或田土核查的本年文书为财政改革史提供节点，制度名称不可跨区套用。（材料核查重点：报称信息的来源、抄传与行政分类）}]
  \item[\eventanchor{ML-1534-001}{1534（嘉靖十三年）：嘉靖倭寇突袭：俘获一名率领50多艘大船的海盗}]
  \item[\eventanchor{ML-1534-002}{1534（嘉靖十三年）：嘉靖皇帝停止例行朝见}]
  \item[\eventanchor{ML-1534-003}{1534（嘉靖十三年）：海禁、朝贡和沿海港口管理的本年材料显示官方海上政策，禁令范围与实际航行须分别调查。（材料核查重点：诏令或奏议的形成与发受链）}]
  \item[\eventanchor{ML-1534-004}{1534（嘉靖十三年）：灾异、赈济及地方工程的本年报告可与州县材料互校，救济命令不自动证明救济到达。}]
  \item[\eventanchor{ML-1534-005}{1534（嘉靖十三年）：海禁、朝贡和沿海港口管理的本年材料显示官方海上政策，禁令范围与实际航行须分别调查。（材料核查重点：报称信息的来源、抄传与行政分类）}]
  \item[\eventanchor{ML-1535-001}{1535（嘉靖十四年）：辽东、广宁守军叛乱，被镇压}]
  \item[\eventanchor{ML-1535-002}{1535（嘉靖十四年）：祀典、学校、出版或讲学的本年材料记录文化制度活动，规范话语不等于社会普遍认同。（材料核查重点：诏令或奏议的形成与发受链）}]
  \item[\eventanchor{ML-1535-003}{1535（嘉靖十四年）：祀典、学校、出版或讲学的本年材料记录文化制度活动，规范话语不等于社会普遍认同。（材料核查重点：报称信息的来源、抄传与行政分类）}]
  \item[\eventanchor{ML-1535-004}{1535（嘉靖十四年）：嘉靖朝礼制、内阁与人事的本年诏令和奏议形成中枢协调线索，留中与施行之间应予区分。}]
  \item[\eventanchor{ML-1536-001}{1536（嘉靖十五年）：蒙古人袭击山西但被击退}]
  \item[\eventanchor{ML-1536-002}{1536（嘉靖十五年）：灾异、赈济及地方工程的本年报告可与州县材料互校，救济命令不自动证明救济到达。（材料核查重点：诏令或奏议的形成与发受链）}]
  \item[\eventanchor{ML-1536-003}{1536（嘉靖十五年）：灾异、赈济及地方工程的本年报告可与州县材料互校，救济命令不自动证明救济到达。（材料核查重点：报称信息的来源、抄传与行政分类）}]
  \item[\eventanchor{ML-1536-004}{1536（嘉靖十五年）：北边警报、边镇防务和西南处置的本年军政材料标记边疆压力，敌我称谓与军数应回到原文。}]
  \item[\eventanchor{ML-1536-005}{1536（嘉靖十五年）：灾异、赈济及地方工程的本年报告可与州县材料互校，救济命令不自动证明救济到达。（材料核查重点：同年地方材料所见的执行范围）}]
  \item[\eventanchor{ML-1537-001}{1537（嘉靖十六年）：蒙古人袭击大同}]
  \item[\eventanchor{ML-1537-002}{1537（嘉靖十六年）：嘉靖朝礼制、内阁与人事的本年诏令和奏议形成中枢协调线索，留中与施行之间应予区分。（材料核查重点：诏令或奏议的形成与发受链）}]
  \item[\eventanchor{ML-1537-003}{1537（嘉靖十六年）：嘉靖朝礼制、内阁与人事的本年诏令和奏议形成中枢协调线索，留中与施行之间应予区分。（材料核查重点：报称信息的来源、抄传与行政分类）}]
  \item[\eventanchor{ML-1537-004}{1537（嘉靖十六年）：赋役折色、盐课、河工或田土核查的本年文书为财政改革史提供节点，制度名称不可跨区套用。}]
  \item[\eventanchor{ML-1537-005}{1537（嘉靖十六年）：嘉靖朝礼制、内阁与人事的本年诏令和奏议形成中枢协调线索，留中与施行之间应予区分。（材料核查重点：同年地方材料所见的执行范围）}]
  \item[\eventanchor{ML-1538-001}{1538（嘉靖十七年）：北边警报、边镇防务和西南处置的本年军政材料标记边疆压力，敌我称谓与军数应回到原文。（材料核查重点：诏令或奏议的形成与发受链）}]
  \item[\eventanchor{ML-1538-002}{1538（嘉靖十七年）：北边警报、边镇防务和西南处置的本年军政材料标记边疆压力，敌我称谓与军数应回到原文。（材料核查重点：报称信息的来源、抄传与行政分类）}]
  \item[\eventanchor{ML-1538-003}{1538（嘉靖十七年）：海禁、朝贡和沿海港口管理的本年材料显示官方海上政策，禁令范围与实际航行须分别调查。}]
  \item[\eventanchor{ML-1538-004}{1538（嘉靖十七年）：北边警报、边镇防务和西南处置的本年军政材料标记边疆压力，敌我称谓与军数应回到原文。（材料核查重点：同年地方材料所见的执行范围）}]
  \item[\eventanchor{MM-1538-GREAT-RITES-INSTITUTIONALIZATION}{1538（嘉靖十七年）：嘉靖进一步调整生父母尊号与祀典，将大礼议的皇室礼制取向制度化}]
  \item[\eventanchor{ML-1539-001}{1539（嘉靖十八年）：日本使团到明中国：日本使节抵达中国后被逮捕并禁止进行贸易}]
  \item[\eventanchor{ML-1539-002}{1539（嘉靖十八年）：辽东守军叛乱，被镇压}]
  \item[\eventanchor{ML-1539-003}{1539（嘉靖十八年）：赋役折色、盐课、河工或田土核查的本年文书为财政改革史提供节点，制度名称不可跨区套用。}]
  \item[\eventanchor{ML-1540-001}{1540（嘉靖十九年）九月：嘉靖皇帝宣布要闭关数年求仙；法院官员说这是无稽之谈并被折磨致死}]
  \item[\eventanchor{ML-1540-002}{1540（嘉靖十九年）：海禁、朝贡和沿海港口管理的本年材料显示官方海上政策，禁令范围与实际航行须分别调查。}]
  \item[\eventanchor{ML-1540-003}{1540（嘉靖十九年）：灾异、赈济及地方工程的本年报告可与州县材料互校，救济命令不自动证明救济到达。}]
  \item[\eventanchor{MM-1540-ALTAN-RAIDS}{1540（嘉靖十九年）：俺答部在大同、宣府方向活动和袭扰，边镇警报上升}]
  \item[\eventanchor{ML-1541-001}{1541（嘉靖二十年）四月30日：一场大火烧毁了太庙大院。}]
  \item[\eventanchor{ML-1541-002}{1541（嘉靖二十年）十月：俺答汗袭陕西。}]
  \item[\eventanchor{ML-1541-003}{1541（嘉靖二十年）：火药在明代用于水利工程。}]
  \item[\eventanchor{ML-1541-004}{1541（嘉靖二十年）：嘉靖朝礼制、内阁与人事的本年诏令和奏议形成中枢协调线索，留中与施行之间应予区分。}]
  \item[\eventanchor{ML-1542-001}{1542（嘉靖二十一年）七月：俺答汗再次袭扰陕西。}]
  \item[\eventanchor{ML-1542-002}{1542（嘉靖二十一年）八月4日：明军在光武之战被阿勒坦汗击败。}]
  \item[\eventanchor{ML-1542-003}{1542（嘉靖二十一年）八月8日：俺答汗掠夺太原郊区。}]
  \item[\eventanchor{ML-1542-004}{1542（嘉靖二十一年）：壬寅宫叛乱：方贵妃阻止嘉靖皇帝被刺杀。}]
  \item[\eventanchor{ML-1542-005}{1542（嘉靖二十一年）：嘉靖皇帝完全退出了正式的职责，在长生宫度过了余生，痴迷于通过药物、仪式和秘传的养生方法来获得身体的不朽。}]
  \item[\eventanchor{MM-1542-RENIN-PALACE-PLOT}{1542（嘉靖二十一年）十月：宫女在乾清宫谋害嘉靖帝未遂，涉案宫女与端妃等随后遭处置}]
  \item[\eventanchor{ML-1543-001}{1543（嘉靖二十二年）十二月：新太庙动工兴建。}]
  \item[\eventanchor{ML-1543-002}{1543（嘉靖二十二年）：浙江饥荒袭来。}]
  \item[\eventanchor{ML-1543-003}{1543（嘉靖二十二年）：祀典、学校、出版或讲学的本年材料记录文化制度活动，规范话语不等于社会普遍认同。}]
  \item[\eventanchor{ML-1544-001}{1544（嘉靖二十三年）：日本对明朝中国的使团：明朝官员拒绝会见日本使节。}]
  \item[\eventanchor{ML-1544-002}{1544（嘉靖二十三年）：浙江又闹饥荒。}]
  \item[\eventanchor{ML-1544-003}{1544（嘉靖二十三年）：灾异、赈济及地方工程的本年报告可与州县材料互校，救济命令不自动证明救济到达。}]
  \item[\eventanchor{ML-1544-004}{1544（嘉靖二十三年）：赋役折色、盐课、河工或田土核查的本年文书为财政改革史提供节点，制度名称不可跨区套用。}]
  \item[\eventanchor{ML-1545-001}{1545（嘉靖二十四年）一月：北京爆发鼠疫。}]
  \item[\eventanchor{ML-1545-002}{1545（嘉靖二十四年）四月：沙尘暴摧毁了冬小麦和大麦作物。}]
  \item[\eventanchor{ML-1545-003}{1545（嘉靖二十四年）七月：新太庙落成。}]
  \item[\eventanchor{ML-1545-004}{1545（嘉靖二十四年）：大同起义，被镇压。}]
  \item[\eventanchor{ML-1545-005}{1545（嘉靖二十四年）：日本使团到明中国：王直随日本使团返回日本，并率领贸易使团前往双榆。}]
  \item[\eventanchor{MM-1545-ALTAN-TRIBUTE-DEBATE}{1545（嘉靖二十四年）：俺答部提出封贡互市要求，明廷围绕受封、贡市与边防反复议论}]
  \item[\eventanchor{ML-1546-001}{1546（嘉靖二十五年）：赋役折色、盐课、河工或田土核查的本年文书为财政改革史提供节点，制度名称不可跨区套用。}]
  \item[\eventanchor{ML-1546-002}{1546（嘉靖二十五年）：嘉靖朝礼制、内阁与人事的本年诏令和奏议形成中枢协调线索，留中与施行之间应予区分。}]
  \item[\eventanchor{ML-1546-003}{1546（嘉靖二十五年）：海禁、朝贡和沿海港口管理的本年材料显示官方海上政策，禁令范围与实际航行须分别调查。}]
  \item[\eventanchor{MM-1546-YAN-SONG-ASCENT}{1546（嘉靖二十五年）：严嵩在内阁中居要位，逐渐取得对奏议传达和人事的影响}]
  \item[\eventanchor{ML-1547-001}{1547（嘉靖二十六年）：嘉靖倭口搜查：监察员报告称，东南沿海的海盗活动已失控。}]
  \item[\eventanchor{ML-1547-002}{1547（嘉靖二十六年）：北边警报、边镇防务和西南处置的本年军政材料标记边疆压力，敌我称谓与军数应回到原文。}]
  \item[\eventanchor{ML-1547-003}{1547（嘉靖二十六年）：祀典、学校、出版或讲学的本年材料记录文化制度活动，规范话语不等于社会普遍认同。}]
  \item[\eventanchor{ML-1548-001}{1548（嘉靖二十七年）二月：嘉靖倭口劫掠：海盗劫掠宁波、台州。}]
  \item[\eventanchor{ML-1548-003}{1548（嘉靖二十七年）六月：蒙古人在宣府击败明军。}]
  \item[\eventanchor{ML-1548-004}{1548（嘉靖二十七年）十月：蒙古人袭击怀来。}]
  \item[\eventanchor{ML-1548-005}{1548（嘉靖二十七年）：明军开始配备火绳枪。}]
  \item[\eventanchor{MM-1548-XIA-YAN-EXECUTION}{1548（嘉靖二十七年）：夏言因边务案件被处死}]
  \item[\eventanchor{ML-1549-001}{1549（嘉靖二十八年）：赋役折色、盐课、河工或田土核查的本年文书为财政改革史提供节点，制度名称不可跨区套用。}]
  \item[\eventanchor{ML-1549-002}{1549（嘉靖二十八年）：灾异、赈济及地方工程的本年报告可与州县材料互校，救济命令不自动证明救济到达。}]
  \item[\eventanchor{MM-1549-DATONG-MUTINY}{1549（嘉靖二十八年）：大同军镇因军饷、军纪等问题发生兵变，朝廷调兵处置}]
  \item[\eventanchor{MM-1549-ZHU-WAN-DEATH}{1549（嘉靖二十八年）：朱纨在遭弹劾、待讯的政治压力下自尽，严厉海禁执行路线受到挫折}]
  \item[\eventanchor{ML-1550-001}{1550（嘉靖二十九年）十月1日：阿勒坦汗掠夺北京郊区}]
  \item[\eventanchor{ML-1550-002}{1550（嘉靖二十九年）十月6日：明军被蒙古人击败}]
  \item[\eventanchor{ML-1550-003}{1550（嘉靖二十九年）：浙江村镇竖起栅栏应对土匪}]
  \item[\eventanchor{ML-1550-004}{1550（嘉靖二十九年）：封贡、互市、月港和港口管理的本年文书记录边海关系重组，官方许可不等于贸易全貌。}]
  \item[\eventanchor{ML-1550-005}{1550（嘉靖二十九年）：本年灾情、赈恤或河工记录连接中央与地方社会，区域差异需以府县材料追踪。}]
  \item[\eventanchor{MM-1550-GENGXU-RAID}{1550（嘉靖二十九年）八月：俺答部抵达北京近郊，京师戒严，史称庚戌之变}]
  \item[\eventanchor{ML-1551-001}{1551（嘉靖三十年）：禁止渔船出海}]
  \item[\eventanchor{ML-1551-002}{1551（嘉靖三十年）：严嵩时期至隆庆初的人事和奏议材料标记中枢权力变化，案狱与党争标签需保留来源立场。（材料核查重点：诏令或奏议的形成与发受链）}]
  \item[\eventanchor{ML-1551-003}{1551（嘉靖三十年）：俺答边患、东南海防或沿海武装的本年军令与战报构成危机处理线索，战果和参与者须交叉核对。（材料核查重点：诏令或奏议的形成与发受链）}]
  \item[\eventanchor{ML-1551-005}{1551（嘉靖三十年）：俺答边患、东南海防或沿海武装的本年军令与战报构成危机处理线索，战果和参与者须交叉核对。（材料核查重点：报称信息的来源、抄传与行政分类）}]
  \item[\eventanchor{MM-1551-ALTAN-SHANXI-RAID}{1551（嘉靖三十年）：俺答部再入山西北部活动，明廷调集边兵并讨论防线与互市}]
  \item[\eventanchor{ML-1552-001}{1552（嘉靖三十一年）四月：明军在大同以北被蒙古人击败}]
  \item[\eventanchor{ML-1552-002}{1552（嘉靖三十一年）：嘉靖倭口突袭：突袭方攻击浙江沿海}]
  \item[\eventanchor{ML-1552-003}{1552（嘉靖三十一年）：嘉靖皇帝选拔800名8岁至14岁的少女入宫}]
  \item[\eventanchor{ML-1552-004}{1552（嘉靖三十一年）：海防知识、学校、科举与礼制的本年文本提供制度和知识史线索，方案不等于实际配置。}]
  \item[\eventanchor{MM-1552-ALTAN-DATONG-RAID}{1552（嘉靖三十一年）：俺答部在大同一线继续入掠，边镇以调兵、守城和转运应对}]
  \item[\eventanchor{ML-1553-001}{1553（嘉靖三十二年）：嘉靖倭寇劫掠：王直劫掠台州以北的浙江沿海}]
  \item[\eventanchor{ML-1553-002}{1553（嘉靖三十二年）：海防知识、学校、科举与礼制的本年文本提供制度和知识史线索，方案不等于实际配置。}]
  \item[\eventanchor{ML-1553-003}{1553（嘉靖三十二年）：军饷、盐课、漕运和赋役的本年行政材料显示中晚明财政压力，法定额与实收不可互代。}]
  \item[\eventanchor{ML-1553-004}{1553（嘉靖三十二年）：本年灾情、赈恤或河工记录连接中央与地方社会，区域差异需以府县材料追踪。}]
  \item[\eventanchor{MM-1553-SHANGHAI-WALL}{1553（嘉靖三十二年）：上海县为应对沿海袭扰修筑城墙，官府、士绅与居民参与筹办}]
  \item[\eventanchor{ML-1554-001}{1554（嘉靖三十三年）春：北京爆发严重疫情}]
  \item[\eventanchor{ML-1554-002}{1554（嘉靖三十三年）三月：嘉靖倭口突袭：海盗杀死松江知县，占领崇明岛}]
  \item[\eventanchor{ML-1554-003}{1554（嘉靖三十三年）：葡中协定（1554年）：莱昂内尔·德索萨贿赂海防副专员，让葡萄牙人留在澳门，每年支付500两银子，并对一半产品征收20\%的关税}]
  \item[\eventanchor{ML-1554-004}{1554（嘉靖三十三年）：俺答边患、东南海防或沿海武装的本年军令与战报构成危机处理线索，战果和参与者须交叉核对。}]
  \item[\eventanchor{MM-1554-WANGJIANGJING}{1554（嘉靖三十三年）：张经、俞大猷等在王江泾击败一支海上武装}]
  \item[\eventanchor{ML-1555-001}{1555（嘉靖三十四年）：嘉靖倭口劫掠：海盗袭击杭州}]
  \item[\eventanchor{ML-1555-003}{1555（嘉靖三十四年）：嘉靖皇帝选拔180名10岁以下的少女入宫}]
  \item[\eventanchor{MM-1555-YANG-JISHENG-EXECUTION}{1555（嘉靖三十四年）：杨继盛在狱中被处死，其弹劾严嵩的奏疏随后广泛流传}]
  \item[\eventanchor{MM-1555-ZHANG-JING-EXECUTION}{1555（嘉靖三十四年）：张经在王江泾获胜后仍因海防案件被处死，东南督师人事骤变}]
  \item[\eventanchor{ML-1556-001}{1556（嘉靖三十五年）一月：1556年陕西大地震：陕西发生地震，死亡人数超过80万}]
  \item[\eventanchor{ML-1556-002}{1556（嘉靖三十五年）：嘉靖倭口劫掠：海盗袭击从南京到杭州的整个海岸线}]
  \item[\eventanchor{ML-1556-003}{1556（嘉靖三十五年）：嘉靖皇帝请求礼部寻找一些神奇的植物来让他长生不老}]
  \item[\eventanchor{ML-1556-004}{1556（嘉靖三十五年）：军饷、盐课、漕运和赋役的本年行政材料显示中晚明财政压力，法定额与实收不可互代。}]
  \item[\eventanchor{MM-1556-SHAANXI-EARTHQUAKE}{1556（嘉靖三十四年十二月）：华县大地震波及关中及邻近地区，朝廷下诏赈恤、修复}]
  \item[\eventanchor{ML-1557-001}{1557（嘉靖三十六年）五月：故宫三大殿毁于火灾}]
  \item[\eventanchor{ML-1557-002}{1557（嘉靖三十六年）冬：阿勒坦汗之子僧格围攻大同附近的驻军}]
  \item[\eventanchor{ML-1557-003}{1557（嘉靖三十六年）：封贡、互市、月港和港口管理的本年文书记录边海关系重组，官方许可不等于贸易全貌。}]
  \item[\eventanchor{MM-1557-MACAU-SETTLEMENT}{1557（嘉靖三十六年）：葡萄牙商人获准在澳门一带停泊、居留和贸易，地方官府对其征收地租等费用}]
  \item[\eventanchor{ML-1558-001}{1558（嘉靖三十七年）四月：嘉靖倭口劫掠：海盗袭击浙江、闽北}]
  \item[\eventanchor{ML-1558-002}{1558（嘉靖三十七年）：援军到来后，僧格撤退}]
  \item[\eventanchor{ML-1558-003}{1558（嘉靖三十七年）：礼部向嘉靖皇帝进贡1860株神奇植物}]
  \item[\eventanchor{ML-1558-004}{1558（嘉靖三十七年）：帝国国库白银跌至不足20万盎司}]
  \item[\eventanchor{MM-1558-QI-JIGUANG-TRAINING}{1558（嘉靖三十七年）：戚继光在浙江募练义乌等地兵员，形成适应沿海作战的部队}]
  \item[\eventanchor{ML-1559-002}{1559（嘉靖三十八年）：长江三角洲地区干旱导致饥荒}]
  \item[\eventanchor{ML-1559-003}{1559（嘉靖三十八年）十二月：苏州卫戍兵变}]
  \item[\eventanchor{ML-1559-004}{1559（嘉靖三十八年）：嘉靖倭口劫掠：海盗占领金门岛并袭击福建和广东}]
  \item[\eventanchor{MM-1559-HU-ZONGXIAN-SEA-DEFENSE}{1559（嘉靖三十八年）：胡宗宪督理浙直军务，以招抚、离间和军事追击并用应对海上武装}]
  \item[\eventanchor{MM-1559-WANG-ZHI-EXECUTION}{1559（嘉靖三十八年）：王直在杭州被处决，明廷以司法处置结束其招抚案件}]
  \item[\eventanchor{ML-1560-001}{1560（嘉靖三十九年）三月：南京卫戍叛军回应削减口粮，直到给他们4万两白银}]
  \item[\eventanchor{ML-1560-002}{1560（嘉靖三十九年）：嘉靖皇帝患有失眠症}]
  \item[\eventanchor{ML-1560-003}{1560（嘉靖三十九年）：戚继光发表了《急效新书》，描述了火枪齐射技术以及他训练明军使用该技术的经验。}]
  \item[\eventanchor{ML-1560-004}{1560（嘉靖三十九年）：封贡、互市、月港和港口管理的本年文书记录边海关系重组，官方许可不等于贸易全貌。}]
  \item[\eventanchor{ML-1560-005}{1560（嘉靖三十九年）：严嵩时期至隆庆初的人事和奏议材料标记中枢权力变化，案狱与党争标签需保留来源立场。}]
  \item[\eventanchor{ML-1561-001}{1561（嘉靖四十年）十二月：紫禁城的住宅宫殿在一场大火中被毁}]
  \item[\eventanchor{ML-1561-002}{1561（嘉靖四十年）：明朝开始生产便携式后装火器。}]
  \item[\eventanchor{ML-1561-003}{1561（嘉靖四十年）：严嵩时期至隆庆初的人事和奏议材料标记中枢权力变化，案狱与党争标签需保留来源立场。}]
  \item[\eventanchor{ML-1561-004}{1561（嘉靖四十年）：海防知识、学校、科举与礼制的本年文本提供制度和知识史线索，方案不等于实际配置。}]
  \item[\eventanchor{MM-1561-TAIZHOU-VICTORIES}{1561（嘉靖四十年）：戚继光部在台州一线连续击败海上武装，巩固浙东防务}]
  \item[\eventanchor{ML-1562-001}{1562（嘉靖四十一年）六月：故宫的住宅宫殿被重建}]
  \item[\eventanchor{ML-1562-002}{1562（嘉靖四十一年）十二月：嘉靖倭口劫掠：海盗占领兴化}]
  \item[\eventanchor{ML-1562-003}{1562（嘉靖四十一年）：海防知识、学校、科举与礼制的本年文本提供制度和知识史线索，方案不等于实际配置。}]
  \item[\eventanchor{MM-1562-FUJIAN-CAMPAIGNS}{1562（嘉靖四十一年）：戚继光等率军进入福建，围攻并清除若干沿海武装据点}]
  \item[\eventanchor{ML-1563-001}{1563（嘉靖四十二年）五月：嘉靖倭寇突袭：明军夺回兴化并摧毁福建海盗基地}]
  \item[\eventanchor{ML-1563-002}{1563（嘉靖四十二年）：海盗林道前遭明军追击，撤退至台湾西南部}]
  \item[\eventanchor{ML-1563-003}{1563（嘉靖四十二年）：明朝将军下令在澎湖修建城墙}]
  \item[\eventanchor{ML-1563-004}{1563（嘉靖四十二年）：俺答边患、东南海防或沿海武装的本年军令与战报构成危机处理线索，战果和参与者须交叉核对。}]
  \item[\eventanchor{MM-1563-PINGHAI-RECOVERY}{1563（嘉靖四十二年）：官军收复平海卫等福建沿海据点}]
  \item[\eventanchor{ML-1564-001}{1564（嘉靖四十三年）：太监把桃子扔到嘉靖皇帝的床上并告诉他这是从天上掉下来的}]
  \item[\eventanchor{ML-1564-002}{1564（嘉靖四十三年）：嘉靖皇帝因应俸禄要求，将所有宗族降为平民}]
  \item[\eventanchor{ML-1564-003}{1564（嘉靖四十三年）：本年灾情、赈恤或河工记录连接中央与地方社会，区域差异需以府县材料追踪。}]
  \item[\eventanchor{ML-1564-004}{1564（嘉靖四十三年）：军饷、盐课、漕运和赋役的本年行政材料显示中晚明财政压力，法定额与实收不可互代。}]
  \item[\eventanchor{MM-1564-QI-JIGUANG-FUJIAN-COMMAND}{1564（嘉靖四十三年）：戚继光继续统领福建陆路军务，与俞大猷等分担沿海据点的清剿和守备}]
  \item[\eventanchor{ML-1565-001}{1565（嘉靖四十四年）：嘉靖皇帝病重}]
  \item[\eventanchor{ML-1565-002}{1565（嘉靖四十四年）：本年灾情、赈恤或河工记录连接中央与地方社会，区域差异需以府县材料追踪。}]
  \item[\eventanchor{ML-1565-003}{1565（嘉靖四十四年）：俺答边患、东南海防或沿海武装的本年军令与战报构成危机处理线索，战果和参与者须交叉核对。}]
  \item[\eventanchor{ML-1565-004}{1565（嘉靖四十四年）：封贡、互市、月港和港口管理的本年文书记录边海关系重组，官方许可不等于贸易全貌。}]
  \item[\eventanchor{MM-1565-HU-ZONGXIAN-ARREST}{1565（嘉靖四十四年）：严嵩失势后，胡宗宪被逮治并死于狱中}]
  \item[\eventanchor{ML-1566-001}{1566（嘉靖四十五年）：嘉靖倭寇劫掠：明军铲除江西、广东海盗}]
  \item[\eventanchor{ML-1566-003}{1566（嘉靖四十五年）：军饷、盐课、漕运和赋役的本年行政材料显示中晚明财政压力，法定额与实收不可互代。}]
  \item[\eventanchor{ML-1566-004}{1566（嘉靖四十五年）：严嵩时期至隆庆初的人事和奏议材料标记中枢权力变化，案狱与党争标签需保留来源立场。}]
  \item[\eventanchor{MM-1566-HAI-RUI-MEMORIAL}{1566（嘉靖四十五年）：海瑞上《治安疏》批评嘉靖政治，随后被下狱}]
  \item[\eventanchor{ML-1567-001}{1567（隆庆元年）一月23日：嘉靖皇帝驾崩}]
  \item[\eventanchor{ML-1567-002}{1567（隆庆元年）二月4日：朱载厚即位隆庆皇帝}]
  \item[\eventanchor{ML-1567-003}{1567（隆庆元年）：海外交易禁令解除}]
  \item[\eventanchor{MM-1567-HAI-RUI-RELEASE}{1567（隆庆元年）正月：隆庆新朝释放因治安疏入狱的海瑞，撤销嘉靖末对其的拘系}]
  \item[\eventanchor{MM-1567-LONGQING-ACCESSION}{1567（隆庆元年）正月：嘉靖帝死后，皇太子朱载垕即位，是为隆庆帝}]
  \item[\eventanchor{MM-1567-ZHANG-JUZHENG-ENTERED-GRAND-SECRETARIAT}{1567（隆庆元年）：张居正入阁，参与隆庆朝边务、财政与人事协调}]
\end{description}
